\section{Appendix B: The Faddeev-Senjanovic formalism}\label{app:faddeev-senjanovic}

In this appendix, we briefly review the Faddeev-Senjanovic formalism, which provides a systematic and rigorous path-integral representation for systems with second-class constraints.

\subsection{B.1 The second-class constraints}\label{subsec:fs-second-class-constraints}

We begin with the Euclidean action of the Proca-Chern-Simons theory after the ADM decomposition:

\begin{align}
S_{E} & =\int \mathrm{d}\tau \int _{\Sigma _{\tau}} \mathrm{d}^{2}x\sqrt{ \sigma }\left(\frac{1}{4}N\tilde{F}_{ab}\tilde{F}^{ab}+\frac{1}{2N}\sigma ^{ab}(\partial _{\tau}A_{a}-D_{a}A_{\tau})(\partial _{\tau}A_{b}-D_{b}A_{\tau})+\frac{\mu ^{2}N}{2}\sigma ^{ab}A_{a}A_{b}+\frac{\mu ^{2}}{2N}A_{\tau}^{2}\right. \\
 & \left.-\frac{ik}{4\pi}\varepsilon ^{ab}(A_{\tau}D_{a}A_{b}-A_{a}(\partial _{\tau}A_{b}-D_{b}A_{\tau}))\right)
\end{align}

The canonical momentum conjugate to $\displaystyle{A_{a}}$ is

\begin{align}
\pi ^{a} & = \frac{\partial \mathcal{L}}{\partial(\partial _{\tau}A_{a})} \\
 & =\frac{\sqrt{ \sigma }}{N}\sigma ^{ab}(\partial _{\tau}A_{b}-D_{b}A_{\tau})-\frac{ik}{4\pi}\epsilon^{ab}A_{b}
\end{align}

while the canonical momentum conjugate to $\displaystyle{A_{\tau}}$ gives the primary constraint

\begin{align}
\chi _{1}=\pi ^{\tau} & \approx 0
\end{align}

The Hamiltonian is then

\begin{align}
H & =\int _{\Sigma}\mathrm{d}^{2}x\left(\pi ^{a}\partial _{\tau}A_{a}-\mathcal{L}\right) \\
 & =\int _{\Sigma}\mathrm{d}^{2}x\sqrt{ \sigma }\left(\frac{N}{2\sigma}\sigma _{ab}\pi ^{a}\pi ^{b}+\frac{\pi ^{a}}{\sqrt{ \sigma }}\left( D_{a}A_{\tau}+\frac{ik}{4\pi} \frac{N}{\sqrt{ \sigma }}\sigma _{ab} \epsilon ^{bc}A_{c} \right)-\frac{1}{4}N\tilde{F}_{ab}\tilde{F}^{ab}\right. \\
 & \left.-\frac{N}{2}\left(\mu ^{2}+\left(\frac{k}{4\pi}\right)^{2}\right)\sigma ^{ab}A_{a}A_{b}-\frac{\mu ^{2}}{2N}A_{\tau}^{2}+\frac{ik}{4\pi}\varepsilon ^{ab}A_{\tau}D_{a}A_{b}\right)
\end{align}

Demanding the preservation of the primary constraint $\displaystyle{\chi _{1}=\pi ^{\tau}\approx 0}$ under time evolution yields the secondary constraint

\begin{align}
\chi _{2} & =\left\{\chi _{1},H\right\}_{\text{P}} \\
 & = -\frac{\delta H}{\delta A_{\tau}} \\
 & =D_{a}\pi ^{a}+\frac{\mu ^{2}\sqrt{ \sigma }}{N}A_{\tau}-\frac{ik}{4\pi}\epsilon ^{ab}D_{a}A_{b}\approx 0
\end{align}

The secondary constraint $\displaystyle{\chi _{2}}$ is then preserved under time evolution without generating any further constraints:

\begin{align}
\left\{\chi _{2},H\right\}_{\text{P}} & \approx 0
\end{align}

Calculating the Poisson bracket between the two constraints gives

\begin{align}
\left\{\chi _{1}(x_{1}),\chi _{2}(x_{2})\right\} & =-\frac{\mu ^{2}\sqrt{ \sigma }}{N}\delta ^{2}(x_{1}-x_{2})\neq 0
\end{align}

This implies that the two constraints $\displaystyle{\chi _{1},\chi _{2}}$ are second class, and the constraint matrix

\begin{align}
C_{\alpha \beta}(x_{1},x_{2}) & =\left\{\chi _{\alpha}(x_{1}),\chi _{\beta}(x_{2})\right\} \\
 & =\begin{pmatrix}
0 & -\frac{\mu ^{2}\sqrt{ \sigma }}{N}\delta ^{2}(x_{1}-x_{2}) \\
\frac{\mu ^{2}\sqrt{ \sigma }}{N}\delta ^{2}(x_{1}-x_{2}) & 0
 \end{pmatrix}
\end{align}

is invertible.

\subsection{B.2 The Faddeev-Senjanovic phase-space path integral}\label{subsec:fs-phase-space-path-integral}

For a Hamiltonian system with second-class constraints $\chi _{\alpha}\approx 0$, the Faddeev-Senjanovic formalism gives the phase-space path integral

\begin{align}
Z & =\int \mathcal{D}q\,\mathcal{D}p\,\prod _{\alpha}\delta [\chi _{\alpha}]\,\det {}^{1/2}\!\big\{\chi _{\alpha},\chi _{\beta}\big\}_{\mathrm{P.B.}}\exp \left(i\int \mathrm{d}t\,[p\dot{q}-H_{c}]\right).
\end{align}

Here the insertion of the delta functionals $\displaystyle{\prod _{\alpha}\delta [\chi _{\alpha}]}$ enforces the second-class constraints at each time slice, while the determinant factor $\displaystyle{\det {}^{1/2}\!\big\{\chi _{\alpha},\chi _{\beta}\big\}_{\mathrm{P.B.}}}$ ensures the correct invariant measure on the constrained phase space. In the present discussion, $\displaystyle{Z}$ denotes a Euclidean phase-space path integral and need not be interpreted as a transition amplitude between two fixed configurations.

Applied to the present model, this becomes

\begin{align}
Z & =\int \mathcal{D}A_{a}\,\mathcal{D}\pi ^{a}\,\mathcal{D}A_{\tau}\,\mathcal{D}\pi ^{\tau}\,\delta [\pi ^{\tau}]\,\delta \left[ D_{a}\pi ^{a}+\frac{\mu ^{2}\sqrt{ \sigma }}{N}A_{\tau}-\frac{ik}{4\pi}\epsilon ^{ab}D_{a}A_{b} \right]\,\det {}^{1/2} C \exp \left(-\int \mathrm{d}\tau\int _{\Sigma _{\tau}}\mathrm{d}^{2}x\,[\pi ^{a}\partial _{\tau}A_{a}-H_{c}]\right)
\end{align}

Integrating out the momentum $\displaystyle{\pi ^{\tau}}$, and noting that the determinant factor $\displaystyle{\det {}^{1/2} C}$ is field independent, we obtain

\begin{align}
Z & =\mathcal{N}\int \mathcal{D}A_{a}\mathcal{D}\pi _{a}\mathcal{D}A_{\tau} \delta\left[ A_{\tau}+\frac{N}{\mu ^{2}\sqrt{ \sigma }}\left(D_{a}\pi ^{a}-\frac{ik}{4\pi}\epsilon ^{ab}D_{a}A_{b}\right)\right] \\
 & \times \exp\left(-\int \mathrm{d}\tau \int _{\Sigma}\mathrm{d}^{2}x\sqrt{ \sigma }\left[\frac{N}{2\sigma}\sigma _{ab}\pi ^{a}\pi ^{b}+\frac{1}{4}N\tilde{F}_{ab}\tilde{F}^{ab}+\frac{N}{2}\left(\mu ^{2}+\frac{k^{2}}{16\pi ^{2}}\right)\sigma ^{ab}A_{a}A_{b}+\frac{\mu ^{2}}{2N}A_{\tau}^{2}-\frac{ik}{4\pi}A_{\tau}\varepsilon ^{ab}D_{a}A_{b}\right]\right)
\end{align}

Here $\displaystyle{\mathcal{N}}$ is a field-independent normalization constant. For later convenience, we denote the delta functional $\displaystyle{\delta\left[ A_{\tau}+\frac{N}{\mu ^{2}\sqrt{ \sigma }}\left(D_{a}\pi ^{a}-\frac{ik}{4\pi}\epsilon ^{ab}D_{a}A_{b}\right)\right]}$ by $\displaystyle{\delta [\chi_{2}]}$, and the exponential factor by $\displaystyle{\exp(-S_{E}[A_{a},A_{\tau};\pi ^{a}])}$. One should keep in mind, however, that $\displaystyle{S_{E}[A_{a},A_{\tau};\pi ^{a}]}$ is not the original Euclidean action, but the Hamiltonian action of the constrained system.

\subsection{B.3 Euclidean correlation functions and the contact term}\label{subsec:fs-euclidean-contact-term}

The Euclidean correlation functions are computed by inserting the corresponding operators into the path integral.

\begin{align}
G_{\tau,\tau}(x_{1},x_{2}) & = \frac{\displaystyle{\int \mathcal{D}A_{a}\mathcal{D}\pi _{a}\mathcal{D}A_{\tau}\delta[\chi _{2}]e^{-S_{E}[A_{a},A_{\tau};\pi ^{a}]}A_{\tau}(x_{1})A_{\tau}(x_{2})}}{\displaystyle{\int\mathcal{D}A_{a}\mathcal{D}\pi _{a}\mathcal{D}A_{\tau}\delta[\chi_{2}]e^{-S_{E}[A_{a},A_{\tau};\pi ^{a}]}}}
\end{align}

We first integrate out $\displaystyle{A_{\tau}}$. The delta functional $\displaystyle{\delta [\chi_{2}]}$ enforces $A_{\tau}$ to coincide with the classical solution

\begin{align}
\mathcal{A}_{\tau}^{\mathrm{cl}} & =-\frac{N}{\mu ^{2}\sqrt{ \sigma }}\left(D_{a}\pi ^{a}-\frac{ik}{4\pi}\epsilon ^{ab}D_{a}A_{b}\right)
\end{align}

Therefore the correlation function can be computed by evaluating the path integral with $A_{\tau}$ fixed to $\mathcal{A}_{\tau}^{\mathrm{cl}}$:

\begin{align}
G_{\tau,\tau}(x_{1},x_{2}) & =\frac{\displaystyle{\int \mathcal{D}A_{a}\mathcal{D}\pi ^{a}e^{-S_{E}[A_{a},A_{\tau}=\mathcal{A}_{\tau}^{\mathrm{cl}};\pi ^{a}]}\mathcal{A}_{\tau}^{\mathrm{cl}}(x_{1})\mathcal{A}_{\tau}^{\mathrm{cl}}(x_{2})}}{\displaystyle{\int \mathcal{D}A_{a}\mathcal{D}\pi ^{a}e^{-S_{E}[A_{a},A_{\tau}=\mathcal{A}_{\tau}^{\mathrm{cl}};\pi ^{a}]}}} \\
 & = \braket{ \mathcal{A}_{\tau}^{\mathrm{cl}}(x_{1})\mathcal{A}_{\tau}^{\mathrm{cl}}(x_{2}) }
\end{align}

This is precisely the canonical correlator, because the second-class constraint has already removed $A_{\tau}$ as an independent integration variable.

\subsection{B.4 Comparison with the configuration-space path integral}\label{subsec:fs-configuration-space-comparison}

The key difference between the Faddeev-Senjanovic formalism and the configuration-space path integral is that in the former the second-class constraint is imposed sharply at the level of the path integral, whereas in the latter one first integrates over all components of $A_{\mu}$ and only then isolates the constraint by completing the square. Consequently, the configuration-space path integral contains fluctuations of $A_{\tau}$ around the classical solution $\mathcal{A}_{\tau}^{\mathrm{cl}}$:

\begin{align}
\eta & :=A_{\tau}-\mathcal{A}_{\tau}^{\mathrm{cl}}
\end{align}

The path integral over $\eta$ that appears in the calculation of $\displaystyle{G_{\tau,\tau;\text{path}}(x_{1},x_{2})}$ is Gaussian, with local propagator

\begin{align}
\braket{ \eta(x_{1})\eta(x_{2}) } & =\frac{N(x_{1})}{\mu ^{2}\sqrt{ \sigma }(x_{1})}\delta ^{3}(x_{1}-x_{2})
\end{align}

Accordingly, the full correlator is the sum of the canonical correlator and this local term:

\begin{align}
G_{\mathrm{path};\tau \tau}(x_{1},x_{2}) & =G_{\tau,\tau}(x_{1},x_{2})+\frac{N(x_{1})}{\mu ^{2}\sqrt{ \sigma }(x_{1})}\delta ^{3}(x_{1}-x_{2})
\end{align}

equivalently,

\begin{align}
G_{\tau,\tau}(x_{1},x_{2}) & =G_{\mathrm{path};\tau \tau}(x_{1},x_{2})-\frac{N(x_{1})}{\mu ^{2}\sqrt{ \sigma }(x_{1})}\delta ^{3}(x_{1}-x_{2}).
\end{align}

Thus the Faddeev-Senjanovic formalism makes the logic especially transparent: it directly produces the reduced phase-space, or equivalently canonical, Euclidean correlator, while the configuration-space path integral produces the covariant Euclidean Green's function. The contact term is exactly the contribution of the independent Gaussian fluctuation that is present in the configuration-space formulation but absent once the second-class constraint is imposed sharply.

For correlation functions involving only the spatial components $A_{a}$, there is no analogous extra local term, because those fields remain genuine dynamical variables in both descriptions. After raising the indices, the relation between the two Euclidean correlation functions is therefore

\begin{align}
G^{\mu _{1},\mu _{2}}(x_{1},x_{2}) & =G_{\mathrm{path}}^{\mu _{1},\mu _{2}}(x_{1},x_{2})-\frac{1}{\mu ^{2}}g^{\mu _{1}\tau }g^{\mu _{2}\tau }\frac{N(x_{1})}{\sqrt{ \sigma }(x_{1})}\delta ^{3}(x_{1}-x_{2}),
\end{align}

which is exactly the contact term found in \hyperref[subsec:explicit-euclidean-correlation-function]{Section~4.3} and rederived in \hyperref[subsec:path-integral-contact-term]{Section~4.5}.

\subsection{B.5 Discussion}\label{subsec:fs-discussion}

The advantage of the Faddeev-Senjanovic formalism is that it makes the role of the second-class constraints completely explicit. Rather than starting from a covariant Gaussian integral over all components of $A_{\mu}$, one begins in phase space and inserts the constraint delta functionals together with the determinant dictated by the Poisson brackets of the constraints. In the present model, this procedure removes $A_{\tau}$ exactly and leaves a reduced path integral over the true dynamical variables.

From this point of view, the contact term is not mysterious. It does not arise from any subtlety in the reduced theory itself; rather, it measures the difference between two inequivalent intermediate representations of the same Euclidean system. The configuration-space path integral retains an auxiliary Gaussian fluctuation of $A_{\tau}$, whereas the Faddeev-Senjanovic path integral eliminates that variable from the outset. The missing self-contraction of this auxiliary fluctuation is precisely the local term.
